\documentclass{controle}
\usepackage{main}

\title{Évaluation n°7 : Limites de suites}
\author{Première Spécialité Mathématiques}
\date{9 Février 2026}

% \printanswers{}

\begin{document}
\titleandname{}
\version{}

\begin{questions}
\question Répondre simplement à chacune des questions suivantes :
\begin{parts}
\part Que signifie qu'une suite admet pour limite $5$ ? 
\part Donner les termes $u_0$, $u_1$, $u_2$ et $u°3$ d'une suite $(u_n)$ convergeant \og sans doute \fg vers $5$.
\part Que signifie qu'une suite diverge ?
\part Donner les termes $v_0$, $v_1$, $v_2$ et $v_3$ d'une suite $(v_n)$ qui diverge \og sans doute \fg.
\part Donner les termes $w_0$, $w_1$, $w_2$ et $w_3$ d'une suite $(w_n)$ qui admet $+\infty$ comme limite.
\end{parts}

\begin{solution}
\begin{parts}
\part Une suite admet $5$ pour limite quand ses termes prennent des valeurs de plus en plus proches de $5$ au fur et à mesure que l'indice $n$ est grand.
\part $u_0 = 6$; $u_1 = 5,5$; $u_2 = 5,25$; $u_3 = 5,125$
\part Une suite divergente est une suite qui ne converge pas, c'est-à-dire qui n'admet pas de limite finie.
\part $v_0 = 1$; $v_1 = -1$; $v_2 = 1$; $v_3 = -1$
\part $w_0 = 1$; $w_1 = 2$; $w_2 = 4$; $w_3 = 8$
\end{parts}
\end{solution}
\end{questions}

\newpage{}
\setcounter{page}{1}
\titleandname{}
\version{}

\begin{questions}
\question Répondre simplement à chacune des questions suivantes :
\begin{parts}
\part Que signifie qu'une suite diverge ?
\part Donner les termes $u_0$, $u_1$, $u_2$ et $u_3$ d'une suite $(u_n)$ qui diverge \og sans doute \fg.
\part Que signifie qu'une suite admet pour limite $-5$ ? 
\part Donner les termes $v_0$, $v_1$, $v_2$ et $v_3$ d'une suite $(v_n)$ convergeant \og sans doute \fg vers $-5$.
\part Que signifie qu'une suite admet $-\infty$ comme limite ?
\end{parts}

\begin{solution}
\begin{parts}
\part Une suite divergente est une suite qui ne converge pas, c'est-à-dire qui n'admet pas de limite finie.
\part $u_0 = 1$; $u_1 = -1$; $u_2 = 1$; $u_3 = -1$
\part Une suite admet $-5$ pour limite quand ses termes prennent des valeurs de plus en plus proches de $-5$ au fur et à mesure que l'indice $n$ est grand.
\part $v_0 = 6$; $v_1 = 5,5$; $v_2 = 5,25$; $v_3 = 5,125$
\part Une suite admet $-\infty$ comme limite quand ses termes prennent des valeurs aussi basses que l'on veut au fur et à mesure que l'indice $n$ est grand.
\end{parts}
\end{solution}
\end{questions}
\end{document}